\markdownRendererDocumentBegin
\markdownRendererWarning{The `hybrid` option has been soft-deprecated.}{The `hybrid` option has been soft-deprecated.}{Consider using one of the following better options for mixing TeX and markdown: `contentBlocks`, `rawAttribute`, `texComments`, `texMathDollars`, `texMathSingleBackslash`, and `texMathDoubleBackslash`. For more information, see the user manual at <https://witiko.github.io/markdown/>.}{Consider using one of the following better options for mixing TeX and markdown: `contentBlocks`, `rawAttribute`, `texComments`, `texMathDollars`, `texMathSingleBackslash`, and `texMathDoubleBackslash`. For more information, see the user manual at <https://witiko.github.io/markdown/>.}\markdownRendererSectionBegin
\markdownRendererHeadingOne{Aviso de derechos de autor y uso}\markdownRendererInterblockSeparator
{}© 2025 Felipe Focil. Todos los derechos reservados.\markdownRendererParagraphSeparator
{}Este libro y todo su contenido —incluyendo texto, diagramas, ejemplos de código, explicaciones teóricas y estructuras organizativas— son propiedad intelectual de su autor.\markdownRendererParagraphSeparator
{}Se autoriza la lectura, copia, impresión y distribución gratuita únicamente con permiso expreso del autor.\markdownRendererSoftLineBreak
{}El autor puede otorgar copias o versiones digitales e impresas de este material de manera libre y sin costo, a personas o grupos que él elija.\markdownRendererParagraphSeparator
{}Queda estrictamente prohibida la venta, redistribución masiva o publicación en medios públicos (impresos o digitales) sin consentimiento previo y por escrito del autor. Cualquier modificación, traducción o adaptación deberá conservar esta nota y reconocer la autoría original.\markdownRendererParagraphSeparator
{}Este material puede compartirse o imprimirse de manera gratuita solo si el autor lo autoriza directamente.\markdownRendererInterblockSeparator
{}\markdownRendererSectionBegin
\markdownRendererHeadingTwo{Referencias}\markdownRendererInterblockSeparator
{}El contenido de este libro se inspira parcialmente en materiales ampliamente reconocidos en la comunidad de programación competitiva, entre ellos:\markdownRendererParagraphSeparator
{}Antti Laaksonen (2020). Guide to Competitive Programming: Learning and Improving Algorithms Through Contests, Third Edition. Springer.\markdownRendererParagraphSeparator
{}CP-Algorithms. https://cp-algorithms.com/\markdownRendererSoftLineBreak
{}— Una referencia abierta y colaborativa mantenida por la comunidad de programación competitiva.\markdownRendererInterblockSeparator
{}
\markdownRendererSectionEnd 
\markdownRendererSectionEnd \markdownRendererSectionBegin
\markdownRendererHeadingOne{Edición y Colaboradores}\markdownRendererInterblockSeparator
{}\markdownRendererSectionBegin
\markdownRendererSectionBegin
\markdownRendererHeadingThree{Autor y Editor Principal}\markdownRendererInterblockSeparator
{}\markdownRendererUlBeginTight
\markdownRendererUlItem Ing. Felipe Rafael Focil Mendoza\markdownRendererUlItemEnd 
\markdownRendererUlEndTight \markdownRendererInterblockSeparator
{}
\markdownRendererSectionEnd \markdownRendererSectionBegin
\markdownRendererHeadingThree{Co-editores y colaboradores}\markdownRendererInterblockSeparator
{}\markdownRendererUlBeginTight
\markdownRendererUlItem Fernanda Avedaño\markdownRendererUlItemEnd 
\markdownRendererUlItem Emiliano Castaños\markdownRendererUlItemEnd 
\markdownRendererUlEndTight \markdownRendererInterblockSeparator
{}
\markdownRendererSectionEnd 
\markdownRendererSectionEnd \markdownRendererSectionBegin
\markdownRendererHeadingTwo{Agradecimientos}\markdownRendererInterblockSeparator
{}A los estudiantes que han utilizado versiones previas del Handbook y ofrecido retroalimentación valiosa, ayudando a refinar su contenido y estructura.\markdownRendererParagraphSeparator
{}Al Mtro. Aldo Tobón, por sentar las bases de este proyecto con su Píldora de Programación, que evolucionaría en el Handbook actual, y por su labor como coach e inspiración constante.\markdownRendererParagraphSeparator
{}Con especial gratitud al Dr. Horacio Jarquín, cuyo acompañamiento, guía y exigencia académica han sido fundamentales para la formación del autor, tanto en la programación competitiva como en el pensamiento analítico. Su compromiso, disciplina y pasión por enseñar han dejado una huella profunda en esta obra y en quienes han aprendido bajo su dirección.\markdownRendererParagraphSeparator
{}A todos los participantes, entrenadores y colegas que continúan impulsando la programación competitiva como un espacio de aprendizaje, colaboración y crecimiento.\markdownRendererParagraphSeparator
{}—Felipe Focil\markdownRendererParagraphSeparator
{}“La suerte solo beneficia a la mente preparada.”\markdownRendererInterblockSeparator
{}
\markdownRendererSectionEnd 
\markdownRendererSectionEnd \markdownRendererSectionBegin
\markdownRendererHeadingOne{Prefacio}\markdownRendererInterblockSeparator
{}El objetivo del Competitive Programming Handbook es ofrecer una recopilación práctica de código y tópicos comunes en concursos de programación, especialmente en el contexto de la ICPC (International Collegiate Programming Contest).\markdownRendererParagraphSeparator
{}A menudo, los estudiantes poseen el conocimiento teórico necesario, pero bajo las condiciones de alta presión típicas de la competencia, la implementación puede tornarse imprecisa o lenta.\markdownRendererParagraphSeparator
{}Este handbook busca servir como una guía de referencia rápida y confiable, reuniendo plantillas, técnicas y algoritmos frecuentemente utilizados por equipos universitarios durante el entrenamiento y las competencias.
\markdownRendererSectionEnd \markdownRendererDocumentEnd